\documentclass[11pt]{article}

\usepackage{amsmath,amssymb}
\usepackage{amsthm}
\usepackage[margin=1in]{geometry}
\usepackage{hyperref}
\usepackage[round,authoryear]{natbib}

\newtheorem{result}{Result}

\title{Partial Tax Recovery on Returns\\
  ($\alpha$-Reclaim under Seller and Buyer Remittance)}
\author{Draft notes\\
  Companion to \texttt{welfare\_with\_returns}}
\date{\today}

\begin{document}
\maketitle

%=============================================================================
\section{Setup}
%=============================================================================

The companion note studies the extreme in which the remitter recovers none of
the unit tax $\tau$ from the government when a unit is returned.  Here the
remitter recovers a fraction $\alpha\in[0,1)$ of $\tau$ on each returned
unit, whether the seller or the buyer remits.  Full reclaim is the boundary
$\alpha=1$; zero reclaim is $\alpha=0$.  Private refunds between buyer and
seller are unchanged from the companion note.  Under seller remittance the
seller refunds the tax-inclusive checkout price $p$.  Under buyer remittance
the seller refunds only the pre-tax list price $p-\tau$, and the buyer
reclaims $\alpha\tau$ from the treasury on a return.

Demand $D(p)$, return rate $r(\cdot)$, handling $h$, and market supply
$S(\mu)$ are as in the companion note.  Equilibrium clears at
$D(p)=S(\mu)$, with $\mu$ the remitter- and $\alpha$-dependent effective
margin derived below.  Write $Q\equiv D(p)$.

%=============================================================================
\section{Effective margins}
%=============================================================================

Under seller remittance the seller remits $\tau$ on every gross sale and
recovers $\alpha\tau$ from the government on each return.  The private refund
remains $p$, so the return rate is still $r(p)$.  Expected net receipt per
gross unit is therefore
\begin{equation}
  \mu^{s}(p,\tau;\alpha)
  = p-\tau-r(p)\bigl(p+h-\alpha\tau\bigr).
  \label{eq:mu-s-alpha}
\end{equation}
The term $-\alpha\tau$ inside the return cost is the treasury reclaim.  At
$\alpha=0$ this is the companion-note margin $p-\tau-r(p)(p+h)$.

Under buyer remittance the seller never remits $\tau$.  The buyer remits on
purchase and recovers $\alpha\tau$ on a return, while the seller refunds only
the pre-tax price $p-\tau$.  The buyer's total recovery on a return is
therefore
\[
  (p-\tau)+\alpha\tau = p-(1-\alpha)\tau,
\]
so the return rate is evaluated at $p-(1-\alpha)\tau$.  The seller's
effective margin is
\begin{equation}
  \mu^{b}(p,\tau;\alpha)
  = (p-\tau)
    -r\bigl(p-(1-\alpha)\tau\bigr)\,(p-\tau+h).
  \label{eq:mu-b-alpha}
\end{equation}
At $\alpha=0$ the argument of $r$ collapses to $p-\tau$ and
\eqref{eq:mu-b-alpha} recovers the companion note.  Clearing conditions are
$D(p)=S\bigl(\mu^{s}(p,\tau;\alpha)\bigr)$ and
$D(p)=S\bigl(\mu^{b}(p,\tau;\alpha)\bigr)$.

%=============================================================================
\section{Tax neutrality}
%=============================================================================

At a common checkout price $p$ and tax $\tau$, the two regimes define the
same equilibrium if and only if $\mu^{s}=\mu^{b}$.  Canceling $p-\tau$ gives
\begin{equation}
  r(p)\bigl(p+h-\alpha\tau\bigr)
  =
  r\bigl(p-(1-\alpha)\tau\bigr)\,(p-\tau+h).
  \label{eq:neutrality-alpha}
\end{equation}
\begin{result}[Partial recovery and tax neutrality]
  Seller and buyer remittance are equivalent if and only if
  \eqref{eq:neutrality-alpha} holds.  When $\alpha=1$, both sides equal
  $r(p)(p-\tau+h)$, so neutrality holds for every $r$.  When $\alpha\in[0,1)$,
  neutrality fails for generic $r$, as in the zero-reclaim companion note.
\end{result}
Full reclaim undoes the remittance wedge on returns.  Any incomplete reclaim
leaves a wedge between the seller's return cost and the buyer's recovery on a
return, so which party remits continues to matter.

%=============================================================================
\section{Pass-through}
%=============================================================================

Production-cost and handling pass-through $\rho_c$ and $\rho_h$ are as in the
companion note and are unchanged by reclaim.  What changes is the tax wedge.
Failure of neutrality requires separate checkout prices $p^{s}$ and $p^{b}$.
With $S_{\tau}=0$,
\begin{equation}
  \rho^{s}_{\tau}
  = \frac{-S'(\mu)\,\mu^{s}_{\tau}}{S'(\mu)\,\mu^{s}_{p}-D'(p^{s})},
  \qquad
  \rho^{b}_{\tau}
  = \frac{-S'(\mu)\,\mu^{b}_{\tau}}{S'(\mu)\,\mu^{b}_{p}-D'(p^{b})}.
  \label{eq:rho-tau-general}
\end{equation}
Differentiating \eqref{eq:mu-s-alpha}--\eqref{eq:mu-b-alpha} at fixed checkout
price, and writing $z\equiv p-(1-\alpha)\tau$,
\begin{align}
  \mu^{s}_{p}
  &= 1-r(p)-r'(p)\bigl(p+h-\alpha\tau\bigr),
  &
  \mu^{s}_{\tau}
  &= -(1-\alpha r(p)),
  \label{eq:mu-s-partials}
  \\
  \mu^{b}_{p}
  &= 1-r(z)-r'(z)\,(p-\tau+h),
  &
  \mu^{b}_{\tau}
  &= -\bigl(1-r(z)\bigr)+r'(z)\,(1-\alpha)\,(p-\tau+h).
  \label{eq:mu-b-partials}
\end{align}
Seller remittance lowers $\mu$ by the unclaimed expected tax $1-\alpha r(p)$.
Buyer remittance lowers the buyer's recovery on a return by only $1-\alpha$,
so the return-rate channel in $\mu^{b}_{\tau}$ is scaled by that factor.

At $\tau=0$, both list-price slopes collapse to the companion object
$\mu_p\equiv 1-r(p)-r'(p)(p+h)$, and
\begin{align}
  \mu^{s}_{\tau}
  &= -(1-\alpha r(p)),
  &
  \mu^{b}_{\tau}
  &= -\bigl(1-r(p)-(1-\alpha)\,r'(p)(p+h)\bigr).
  \label{eq:mu-tau-alpha}
\end{align}
Comparing with $\rho_c$ then yields
\begin{equation}
  \rho^{s}_{\tau}
  = (1-\alpha r(p))\,\rho_c,
  \qquad
  \rho^{b}_{\tau}
  = \bigl(1-r(p)-(1-\alpha)\,r'(p)(p+h)\bigr)\,\rho_c.
  \label{eq:rho-alpha}
\end{equation}
\begin{result}[Tax pass-through under partial reclaim]
\label{res:rho-alpha}
  At $\tau=0$, tax pass-through is \eqref{eq:rho-alpha}.  The formulas nest the
  companion note at $\alpha=0$, where $\rho^{s}_{\tau}=\rho_c$ and
  $\rho^{b}_{\tau}=\mu_p\rho_c$, and coincide at $(1-r(p))\rho_c$ when
  $\alpha=1$.  For $\alpha\in[0,1)$,
  $\operatorname{sgn}(\rho^{b}_{\tau})=-\operatorname{sgn}(\rho^{s}_{\tau})$
  if and only if $1-r(p)-(1-\alpha)\,r'(p)(p+h)<0$ and $\rho_c\neq 0$.
\end{result}
At $\tau=0$, the list-price slope $\mu_p$ does not depend on reclaim, so
demand and how supply responds to $p$ are the same as for a production-cost
shock.  Seller-tax pass-through therefore scales only with how much $\tau$
lowers the effective margin---and thus creates excess demand---at a fixed
checkout price.  That shift is $1-\alpha r(p)$: the firm remits one dollar on
every gross sale but recovers $\alpha$ on the fraction $r(p)$ that returns, so
expected unclaimed tax falls in $\alpha$.  Hence
$\rho^{s}_{\tau}=(1-\alpha r)\,\rho_c$, which declines from $\rho_c$ at
$\alpha=0$ toward $(1-r)\rho_c$ as reclaim becomes complete: more complete
recovery leaves a smaller excess-demand gap, so checkout-price pass-through is
weaker.  Cost and handling pass-through are untouched because those shocks
never enter the $\alpha\tau$ reclaim term.

Under buyer remittance, pass-through again scales only with how $\tau$
shifts the effective margin at a fixed checkout price.  The buyer's recovery
on a return is the seller's pre-tax refund plus the treasury reclaim,
$p-(1-\alpha)\tau$.  A one-dollar tax therefore lowers that recovery by
$1-\alpha$, so the return-rate response is scaled by the unreclaimed share.
The resulting factor is $1-r-(1-\alpha)r'(p)(p+h)$.
If $r'<0$, the unreclaimed tax raises the return rate and enlarges the hit to
the effective margin above $1-r$.  More complete reclaim shrinks that channel,
so the factor falls toward $1-r$ and checkout-price pass-through becomes
weaker.
If $r'>0$, the unreclaimed tax lowers the return rate and shrinks the hit to
the effective margin below $1-r$.  More complete reclaim shrinks that channel,
so the factor rises toward $1-r$.

The companion note also signs the classical bias for $\rho_c$: the no-return
formula overestimates pass-through if and only if $\mu_p>1$.  That comparison
does not involve $\alpha$.  At $\tau=0$ the list-price slope is
$\mu_p=1-r(p)-r'(p)(p+h)$ under either remittance regime, because reclaim
enters only through terms in $\alpha\tau$ that drop out of
$\partial\mu/\partial p$.  Production-cost pass-through is
$\rho_c=S'(\mu)/(S'(\mu)\mu_p-D'(p))$, again with no $\alpha$.  The classical
formula sets $\mu_p=1$ in that expression, so whether it over- or
underestimates $\rho_c$ is pinned entirely by whether $\mu_p\gtrless 1$---the
same $\varepsilon_r(p)<-p/(p+h)$ cut as in the companion note, for every
$\alpha\in[0,1]$.  Full reclaim is no exception: equalizing the tax rates
leaves $\rho_c$ and its bias untouched.
\begin{result}[Pass-through bias under partial reclaim]
\label{res:pt-bias-alpha}
  The companion bias cut for $\rho_c$ relative to the classical no-return
  formula is unchanged by $\alpha$, including at full reclaim $\alpha=1$.
  Tax pass-through inherits that bias only through the reclaim factors in
  \eqref{eq:rho-alpha}.
\end{result}

%=============================================================================
\section{Revenue and first-order surplus}
%=============================================================================

On each gross sale the treasury collects $\tau$.  On each return it refunds
$\alpha\tau$ to the remitter.  Net revenue is therefore
\begin{equation}
  G
  = \tau\,Q\bigl(1-\alpha\,r^{\ast}\bigr),
  \label{eq:G-alpha}
\end{equation}
where $r^{\ast}=r(p)$ under seller remittance and
$r^{\ast}=r\bigl(p-(1-\alpha)\tau\bigr)$ under buyer remittance.  At
$\tau=0$ both give $dG/d\tau=Q\bigl(1-\alpha r(p)\bigr)$.

Private surplus $W=CS+PS$ is differentiated as in the companion note, using
$dCS/d\tau=-Q(1-r)\rho_{\tau}$ and
$dPS/d\tau=Q(\mu_p\rho_{\tau}+\mu_{\tau})$ at $\tau=0$.  Substituting
\eqref{eq:mu-tau-alpha}--\eqref{eq:rho-alpha} and writing
$\Lambda\equiv 1+r'(p)(p+h)\rho_c$ yields
\begin{align}
  \frac{dW}{d\tau^{s}}
  &= -Q\bigl(1-\alpha r(p)\bigr)\,\Lambda,
  &
  \frac{dW}{d\tau^{b}}
  &= -Q\bigl(1-r(p)-(1-\alpha)\,r'(p)(p+h)\bigr)\,\Lambda.
  \label{eq:dW-alpha}
\end{align}
Social surplus $V=W+G$ then satisfies, at $\tau=0$,
\begin{align}
  \frac{dV}{d\tau^{s}}
  &= -Q\bigl(1-\alpha r(p)\bigr)\,(\Lambda-1),
  &
  \frac{dV}{d\tau^{b}}
  &= Q\Bigl[
    \bigl(1-\alpha r(p)\bigr)
    -\bigl(1-r(p)-(1-\alpha)\,r'(p)(p+h)\bigr)\,\Lambda
  \Bigr].
  \label{eq:dV-alpha}
\end{align}
When $\alpha=0$, \eqref{eq:dW-alpha}--\eqref{eq:dV-alpha} reproduce the
companion note.  When $\alpha\to 1$, seller remittance has
$dW/d\tau^{s}\to -Q(1-r)\Lambda$ and
$dV/d\tau^{s}\to -Q(1-r)(\Lambda-1)$, so the unclaimed tax on returned units
vanishes from the first-order social wedge and only the kept-unit share
$1-r$ remains in the private loss.  Incomplete reclaim $\alpha<1$ leaves a
strictly larger seller-side resource cost $1-\alpha r>1-r$ whenever $r\in(0,1)$.

\begin{result}[Partial recovery nests zero and full reclaim]
  For $\alpha\in[0,1]$, equations
  \eqref{eq:mu-s-alpha}--\eqref{eq:dV-alpha} nest the zero-reclaim companion
  note at $\alpha=0$.  At $\alpha=1$, tax neutrality is restored and the
  seller-tax pass-through factor collapses to the keep rate $1-r(p)$.
\end{result}

\end{document}
